% Overleaf: main document main.tex; compiler pdfLaTeX.
% The table is embedded. Upload the accompanying figs/ folder.
% tables/tab_profiles.tex is an unchanged reference copy.

% Overleaf: upload this file together with the figures in figs/ (fig1_scopeP_vs_scopeF.png,
% fig2_position_profiles.png, fig2b_query_fact_tokens.png, fig3_final_attention.png,
% fig4_query_control.png, fig5_contextual_output.png, fig6_copying_weights.png,
% fig7_synthetic.png, fig8_attention_value.png, fig9_attended_name.png) and tables/tab_profiles.tex,
% then compile with pdfLaTeX. Source data: results/stage3_v1/analysis/ (run stage3_v1, job completed
% 22 September 2026; summary.json complete=true).
\documentclass[11pt,a4paper]{article}
\usepackage[margin=2.2cm]{geometry}
\usepackage[T1]{fontenc}
\usepackage[utf8]{inputenc}
\IfFileExists{lmodern.sty}{\usepackage{lmodern,microtype}}{}
\usepackage{amsmath,booktabs,graphicx}
\usepackage{xcolor}
\usepackage[colorlinks=true,urlcolor=blue!60!black,linkcolor=blue!60!black]{hyperref}
\graphicspath{{figs/}{./}}
\setlength{\parindent}{0pt}
\setlength{\parskip}{5pt}
\setlength{\emergencystretch}{2em}
\renewcommand{\arraystretch}{1.15}

\title{Stage 3: Results and Analysis\\\large Head-Level Characterization of the Single-hop Causal Heads in OLMo-2-7B}
\author{Daniella Simonovsky and Max Golubkov}
\date{September 2026}

\begin{document}
\maketitle

\section{Run and validation}
Run \texttt{stage3\_v1} completed on the cluster with three model replicas (six GPUs)
and exit code 0. Every replica passed the gate before measurements began: behavioral
replication drift 0, self-patch drift 0, zero change when corrupted activations are
patched before the first differing token, reproduction of the saved Stage-2
\texttt{exact\_delta} on the six probe heads, first-divergent-token alignment,
QK/RoPE-aware reconstruction of the SDPA attention patterns, replication of the saved
raw RI scores within the frozen tolerances, and agreement between the joint
attention+value replacement and the ordinary head-output patch. The gate measured
0.151--0.152~s per patched forward and predicted 3.7~GPU-hours for the 263{,}328
new patched forwards (8{,}188 missing Scope-F patches, 176{,}820 position-scan
patches, 74{,}760 attention/value factorial patches and 3{,}560 reverse patches).

The automatic CPU audit verified every checkpoint checksum, the coverage of all 105
heads on all 178 pairs, the position scan on the 40 common pairs, the 160 core prompts
(20 common families $\times$ two orders $\times$ \{base, corrupted, reorder,
changed-query\}), the 34{,}526 matched contextual-RI anchor comparisons, the weight
diagnostic and both synthetic benchmarks, and wrote \texttt{summary.json} with
\texttt{complete: true}. The 46 heads that lacked a Scope-F measurement now have one
on all 178 pairs; no saved Stage-1/2 value was recomputed or replaced. All results are
descriptive. Causal measurements (Scope P, Scope F, reverse patches, the attention--value
factorial) are family-level means over the 89 discovery families with the two orders
averaged within family; the position scan, attention and output profiles and the
attended-name classification use the 20 common families (the latter counts eligible
events, not families); the synthetic fingerprints average over trials. The 87
validation families remain untouched.

\textbf{What Stage 3 measured.} For the fixed inventory of 105 heads (48 RI-selected
small-effect, 12 strong positive outside RI, 5 strong positive RI-selected, 8 strong
negative, 7 moderate, 25 random controls) we report: (i) Scope-P vs.\ Scope-F
importance with the clean/corrupted logit decomposition; (ii) the causal position
profile (single-position patches at every test-block position, 40 common pairs) and
reverse patching for the ten largest heads; (iii) the unconditional attention profile
at the final colon and the changed-query, reorder and corruption controls; (iv) the
contextual output projection $o^h_j W_U$ against the raw-embedding projection used by
RI; and the four planned diagnostics: (v) weight copying scores, (vi) attended-name
mover classification, (vii) matched contextual RI, (viii) the attention--value
factorial, and (ix) two synthetic fingerprints. Section~\ref{sec:cards} gives the
head-profile table; Section~\ref{sec:hyp} the evidence assessment of the three Stage-2
hypotheses; Section~\ref{sec:stage4} the hand-off to Stage~4.

\section{Notation}
As in Stage 2, $M(x)=\operatorname{logit}(y_c)-\operatorname{logit}(y_r)$ at the first
divergent answer token, with the sign fixed to the clean answer. For head $h$,
$I_h=-[M(x_c;a_h\leftarrow a^r_h)-M(x_c)]$ with the corrupted donor inserted at all
original prompt positions (Scope~P) or at the final colon only (Scope~F, $I^F_h$).
Positive importance means the patch moves the metric toward the corrupted answer. For
every Scope-F patch the run stored the clean- and corrupted-answer logits, so
$I^F_h=P_h+S_h$ where $P_h$ is the rise of the corrupted-answer logit (``promotion'')
and $S_h$ the fall of the clean-answer logit (``suppression''). Reverse recovery is
$M(x_r;a_h\leftarrow a^c_h)-M(x_r)$ (clean donor into the corrupted prompt). The test
block reads ``$\langle$mother$\rangle$ is the mother of $\langle$child$\rangle$.''
four times, then ``Question: Who is the mother of $\langle$child$\rangle$? Answer:''.
The corruption swaps the two answer-side mothers; the \emph{query fact} is the fact
whose child is asked about, and its mother is the answer.

\section{Where the effect enters: Scope P vs.\ Scope F}
\label{sec:scope}
\begin{figure}[ht]
\centering
\includegraphics[width=.72\textwidth]{fig1_scopeP_vs_scopeF.png}
\caption{Scope-P against Scope-F importance for all 105 inventory heads (178 pairs,
symlog axes). Heads on the diagonal deliver their whole effect at the answer position;
heads on the horizontal axis act only at earlier positions.}
\label{fig:scope}
\end{figure}

Figure~\ref{fig:scope} completes the picture that Stage 2 could only sketch for 69
heads. Of the 25 heads with $|I_h|\ge 0.3$, 20 have $I^F_h/I_h$ between 0.96 and 1.07:
the final colon carries their entire effect, and the sign is consistent across
families ($F$ positive in $\ge 97.8\%$ of families for every strong positive head,
in $\le 4.5\%$ for every strong negative head). Five heads are different:

\begin{center}
\small
\begin{tabular}{lrrrl}
\toprule
Head & $I_h$ (P) & $I^F_h$ & $F/P$ & Position of the effect (Section~\ref{sec:position}) \\
\midrule
L17H1  & $+5.701$ & $-0.056$ & $-0.01$ & query-fact sentence: mother tokens, ``is'', period \\
L15H25 & $+0.679$ & $+0.169$ & $0.25$ & last token of the query-fact child \\
L14H26 & $+0.399$ & $+0.007$ & $0.02$ & diffuse; template words \\
L13H10 & $-0.681$ & $-0.029$ & $0.04$ & query-fact sentence: ``is'', ``of'', mother last token \\
L17H3  & $-0.546$ & $-0.402$ & $0.74$ & mixed: colon plus earlier positions \\
\bottomrule
\end{tabular}
\end{center}

The most important single head of the whole project, L17H1 (5.7 logits, half the
clean--corrupted gap), has \emph{no} effect at the answer position. Stage 2 had labelled
the layer-16--26 heads ``readers at the colon'' from 69 heads; L17H1 was not among
them, and the label does not hold for it. Twenty-five random controls have
$|I^F_h|\le 0.019$ except L18H22 ($+0.089$), which also has a Scope-P effect of the
same size.

\textbf{Logit decomposition.} For every strong head the Scope-F effect is dominated by
the corrupted-answer logit: $P_h$ accounts for 93\% (L18H19: $+2.88$ of $+3.08$),
96\% (L18H18), 88\% (L21H18), 93\% (L22H5) and 98\% (L20H1: $-1.49$ of $-1.53$;
L23H15: $-0.94$ of $-0.95$; L19H16: $-0.83$ of $-0.83$). L27H6 is the exception:
$P=+1.96$, $S=+1.10$; it both promotes the donor's name and lowers the other. The
negative heads therefore act by \emph{lowering} the corrupted-answer logit when their
corrupted-run activation is inserted, with almost no change of the clean-answer logit.
Under the sign convention of Stage 2, a head that suppresses the name it retrieves on
its own run produces exactly this pattern (the corrupted-run activation suppresses
the corrupted answer); Section~\ref{sec:attended} tests that reading directly.

\textbf{Reverse patching.} For the ten fixed heads, inserting the clean activation into
the corrupted prompt recovers the metric by almost exactly the noising effect
(L17H1: $+5.21$ vs.\ $5.70$; L18H19 $+3.22$ vs.\ $3.07$; L27H6 $+3.06$; L18H18
$+2.85$; L21H18 $+1.75$; L25H17 $+1.27$; L22H5 $+1.01$; L20H1 $-1.63$; L30H18
$-1.20$; L23H15 $-0.93$). Recovery is symmetric with noising for every head, at
Scope~P and at Scope~F alike (L17H1: $-0.055$ at the colon). This is conditional on
the rest of the corrupted model and does not show that any head suffices alone.

\section{Causal position profile}
\label{sec:position}
\begin{figure}[ht]
\centering
\includegraphics[width=\textwidth]{fig2_position_profiles.png}
\caption{Single-position patches on the common 40 pairs, summed by position role within
each prompt and averaged over the 20 families. Top: the four early-position heads and
two moderate heads; bottom: four answer-position heads. Titles give Scope P and F on the
same 40 pairs.}
\label{fig:position}
\end{figure}
\begin{figure}[ht]
\centering
\includegraphics[width=\textwidth]{fig2b_query_fact_tokens.png}
\caption{Token-by-token single-position importance inside the query-fact sentence
(family mean $\pm$ SD, 20 families) for L17H1, L13H10 and L15H25.}
\label{fig:tokens}
\end{figure}

\textbf{Sums of single-position effects vs.\ the all-position effect.} Summing a
head's single-position effects over all test positions tracks its all-position Scope-P
effect on the same 40 pairs in the aggregate: Spearman 0.913 over the 105 heads, and a
signed mean residual of $-0.135$ logits for L17H1 (5.61 vs.\ 5.74). The signed mean
hides large cancelling errors, however: for L17H1 the mean absolute residual per pair is
0.94 logits (per family, orders averaged, 0.89) and the maximum is 2.85 (2.33 per
family). For the answer-position heads the residuals are small at every level (L18H19:
mean absolute family residual 0.06, maximum 0.14; L27H6: 0.03, 0.09), for the trivial
reason that all of their effect sits at one position. The single-position profile
therefore locates \emph{where} a head's activity influences the answer; it is not a
proven additive partition of the total effect, and within-head interactions between
positions are not established to be small for L17H1.

\textbf{L17H1 acts inside the query-fact sentence, after the mother's first token.}
Figure~\ref{fig:tokens} shows the profile: mother first token $+0.05$, middle
$+1.15$, last $+1.67$, ``is'' $+1.77$, ``the'' $0.00$, ``mother'' $+0.33$, ``of''
$+0.36$, child tokens $\approx 0$, period $+1.38$. Nothing at the question, nothing at
the colon. At the mother tokens the clean and corrupted \emph{inputs} differ, so a
head that writes a representation of the current name would show an effect there for
trivial reasons; the effect at ``is'', ``of'' and the period, whose inputs are
identical in the two prompts, shows that L17H1's output at those positions depends on
the preceding mother and matters for the answer. The natural reading is that it carries
mother-dependent information forward through the sentence into positions later heads
read from; what that information is (the name, the relation, or a binding of the two)
is not identified by this measurement (see Section~\ref{sec:cards} on the optional
readout). The same pattern with
opposite sign describes L13H10 (``is'' $-0.22$, ``of'' $-0.22$, mother last $-0.14$).
At the distractor facts L17H1 has small \emph{negative} effects (distractor mother
$-0.53$ summed), as expected if the corrupted donor writes the wrong name identity into
the other candidate fact.

\textbf{L15H25 acts at the query child.} Its effect is concentrated on the last token
of the query-fact child ($+0.43$ of 0.66), with $+0.13$ at the colon. The child's
tokens are identical in the two prompts; the head therefore writes, at the child's
position, something that depends on the preceding mother. L14H26 (0.40) is diffuse
over template words with no dominant site.

\textbf{Answer-position heads.} L18H19, L27H6, L18H18, L21H18, L25H17, L22H5, L16H21,
L16H1, L21H6, L30H13, L21H23, L17H17, L17H24, L19H22 and all eight strong negative
heads except L13H10 (and partly L17H3, $F/P=0.74$) have flat profiles with a single
bar at the colon.

\section{Relation-tracking attention profile}
\label{sec:attention}
\begin{figure}[ht]
\centering
\includegraphics[width=\textwidth]{fig3_final_attention.png}
\caption{Attention mass at the final colon on base prompts (20 common families, both
orders), for the 22 heads with $|I^F_h|\ge0.3$.}
\label{fig:attention}
\end{figure}
\begin{figure}[ht]
\centering
\includegraphics[width=\textwidth]{fig4_query_control.png}
\caption{Left: paired attention shift when the question is changed to the other
answer-side child (facts unchanged), toward the newly queried mother ($x$) and child
($y$), family means. Right: the same shift beside the paired change under reordering
and under corruption for the strong heads.}
\label{fig:query}
\end{figure}

The answer-position heads fall into three attention types (Figure~\ref{fig:attention};
random controls: mean mass 0.017 on the query mother, maximum 0.109, 0.009 on the
query child).

\emph{Mother-attenders.} L27H6 (0.58 on the query mother), L21H18 (0.60), L23H15
(0.51), L22H5 (0.34), L19H22 (0.31), L21H6 (0.30), L19H16 (0.21), L26H23 (0.21),
L20H1 (0.17), L18H18 (0.14), L18H19 (0.13), L17H24 (0.13). They attend to the answer's
own mention. Both signs occur: L23H15, L19H16, L26H23 and L20H1 attend to the same
place as L27H6 and L21H18 but have negative effects.

\emph{Child-attenders.} L16H21 (0.40 on the query child, 0.01 on the mother), L16H1
(0.36), L16H24 (0.37, small effect), L16H31 (0.14, negative). They attend to the
second mention of the queried child inside the facts, not to the answer.

\emph{Self-attenders.} L25H17 (0.57 on the colon itself, 1.27 logits), L13H10 (0.42),
L21H23 (0.20), L30H13 (0.19), L17H17 (0.19). Self attention is their largest single
target, but it is not exclusive: the eight fact-name spans together receive 0.03
(L25H17, L13H10), 0.09 (L30H13), 0.15 (L21H23) and 0.22 (L17H17) of their mass, and the
remainder is spread over other positions. For L25H17 most of the value it carries
comes from the colon's own residual stream; for L17H17 and L21H23 a fact-reading path
cannot be excluded.

L30H18 ($-1.06$) is the one strong head that attends more to the \emph{distractor}
mothers (0.11 per fact) than to the query mother (0.05).

\textbf{Attention follows the question, not the location or the token
(Figure~\ref{fig:query}).} Changing only the queried child moves the mother-attenders'
mass to the newly relevant fact in 100\% of families (L21H18 $+1.07$, L23H15 $+0.91$,
L27H6 $+0.84$, L19H22 $+0.56$, L21H6 $+0.50$, L22H5 $+0.49$, L26H23 $+0.37$, L19H16
$+0.33$, L20H1 $+0.32$, L18H19 $+0.23$, L18H18 $+0.20$), and moves the
child-attenders' mass to the new child (L16H21 $+0.83$, L16H1 $+0.68$, L16H31
$+0.27$). Random controls: at most $+0.18$. (The preference shift is a difference of
two attention contrasts, so values above 1 are possible.) Reordering the facts and
corrupting the mother change the family-mean attention on the query mother by at most
0.11 (L21H6) and 0.085 (L27H6) respectively, but these signed means also hide
prompt-level variation: for L21H18 the mean absolute change under reordering is 0.15
per prompt (maximum 0.49) and under corruption 0.12; for L27H6 0.09 and 0.29; for
L23H15 0.12 and 0.14 (the child-attender L16H21: 0.004 and 0.004). The heads therefore
keep attending to the relevant fact \emph{on average} when it is moved or its name is
swapped, but positional invariance at the level of individual prompts is not
established; the changed-query result is the stronger evidence, because it is
consistent in sign in every family. The model itself answers the changed query correctly in 38 of the 40 changed-query
prompts (mean margin 5.25 vs.\ 5.85 on the base prompts). L30H18 is again the
exception: its mass moves \emph{away} from the newly queried fact ($-0.19$, negative
in every family). None of this establishes what the heads compute; it establishes that
their routing is query-selective on these 20 families.

\section{Contextual output profile}
\label{sec:output}
\begin{figure}[ht]
\centering
\includegraphics[width=\textwidth]{fig5_contextual_output.png}
\caption{Scope-F importance against the vocabulary projection of the head's output at
the colon (query-mother first token minus mean distractor-mother first token). Left:
actual contextual output $o^h_j W_U$. Right: the raw-embedding projection of the colon
token used by RI, on the same $x$ range as the left panel would need ($\pm0.02$).}
\label{fig:output}
\end{figure}

The raw-embedding projection is uninformative at the colon for every head (all values
within $\pm0.008$; Spearman with $I^F_h$ $-0.07$), as the algebraic argument of Stage 1
predicted. The contextual projection is informative but partial (Spearman $+0.32$).
Three groups emerge.

\emph{Direct writers of the answer.} L27H6 writes $+2.99$ logits for the query mother
against the distractor mothers (first token; at the last-token anchor only $+0.06$); L30H13 $+0.76$; L22H5
$+0.41$; L21H6 $+0.33$; L21H18 $+0.25$. For L27H6 the projection is of the same size
as its causal effect, but the two numbers are not an attribution equality: the
projection is one clean vector through $W_U$ before the shared post-attention norm, on
20 families, contrasting the query mother with the distractor mothers, whereas $I^F_h$
is a clean-to-donor change in the full model on 89 families, at the first divergent
answer token. L30H18 writes $-0.88$: its output lowers the query mother relative to the
distractors, consistent with its distractor-directed attention and its causal sign.

\emph{Heads whose projection does not explain their effect.} L18H19 ($I^F=3.08$)
writes only $+0.16$ for the answer, L18H18 ($2.72$) $+0.11$, L20H1 ($-1.53$) $-0.09$,
L23H15 ($-0.95$) $+0.04$, L19H16 ($-0.83$) $-0.04$, L25H17 ($1.27$) $+0.16$; L16H21 and
L16H1 have contrasts within $\pm0.02$ on both mother and child tokens against
$0.9$-logit effects. For these heads the selected vocabulary projection of the clean
output vector does not account for the measured intervention effect. The hypothesis
this suggests is mediation by later components at the colon (later layers, since the
colon is the last prompt token), but the present data do not establish it: the
projection ignores the shared post-attention normalization, whose response to the
donor vector is itself a candidate explanation, and it is a property of one clean
vector rather than of the clean-minus-donor difference that the intervention inserts.
A matched clean-minus-donor projection under an explicit normalization reference would
be a sharper diagnostic; identifying a downstream path requires the Stage-4 path
interventions.

All values are pre-normalization diagnostics of one head's vector; they are not logit
attributions, and the shared post-attention norm of OLMo-2 is not assigned to heads.

\section{Planned diagnostics}
\subsection{Copying and token preference (weights)}
\begin{figure}[ht]
\centering
\includegraphics[width=\textwidth]{fig6_copying_weights.png}
\caption{Left: weight copying score on the 314 first/last test-name token IDs
($e_tW^h_{OV}W_U$, self minus mean over other names) against signed Scope-P
importance. Right: the same score on names against the 128-token ordinary-word
reference sample.}
\label{fig:copy}
\end{figure}

The static copying score separates the sign of the effect but not its size: Spearman
with signed $I_h$ is $+0.37$ ($n=105$), with $|I_h|$ $+0.02$. Positive copiers with
positive effects: L21H6 ($+0.35$; the token is its own top output for 60\% of names),
L17H24 ($+0.19$), L19H22 ($+0.17$), L18H19 ($+0.13$), L21H18 ($+0.12$). Anti-copiers
with negative effects: L26H23 ($-0.20$), L20H1 ($-0.17$), L19H16 ($-0.09$). Random
controls lie in $[-0.044,+0.044]$. But the two largest heads at the colon are not
copiers by weight (L27H6 $+0.04$, L18H18 $-0.08$), L17H1 is neutral ($+0.01$), and the
strongest name-copier in the inventory, L30H18 ($+0.60$, far outside every other head),
has a \emph{negative} effect because it applies its copying to the distractor it
attends to. Name-token and ordinary-word scores are similar for most heads
(Figure~\ref{fig:copy}, right); L30H18 alone is name-specific ($+0.60$ vs.\ $+0.03$).
Weight alignment therefore neither identifies the important heads nor decides the
copy-penalizing hypothesis on its own; the event-based test follows.


\subsection{Copying and token preference (saved RI events by current-token role)}
\label{sec:current}
The plan's first copying diagnostic partitions the saved Stage-1 test-only RI events by
whether the \emph{current} token belongs to the target name, to a control name, or to
neither, and reports the raw target-minus-control gap in each partition. It was computed
during preparation (\texttt{current\_name\_diagnostics.json}) and is reported here; it
is distinct from the attended-name test, which concerns the token a head attends to.
\begin{table}[ht]
\centering\small
\begin{tabular}{lrlll}
\toprule
Head & $I_h$ & Current token $\in$ target & Current token $\in$ control name & Neither \\
\midrule
L18H19 & +3.07 & $+0.1011$ (1/1)$^{\dagger}$ & $-0.0228$ (111/48) & $+0.0003$ (51/32) \\
L21H18 & +1.70 & -- & $-0.0168$ (194/65) & $+0.0011$ (185/74) \\
L21H6 & +0.70 & -- & $-0.0245$ (421/85) & $-0.0044$ (116/70) \\
L22H5 & +0.98 & -- & $-0.0187$ (10/8)$^{\dagger}$ & $+0.0333$ (14/5)$^{\dagger}$ \\
L19H22 & +0.43 & -- & $-0.0148$ (19/13)$^{\dagger}$ & $-0.0001$ (79/47) \\
L16H21 & +0.92 & -- & $+0.0042$ (149/67) & $-0.0007$ (307/82) \\
L16H1 & +0.87 & $-0.0039$ (88/34) & $-0.0020$ (242/75) & $-0.0041$ (331/82) \\
L25H17 & +1.27 & -- & $+0.0009$ (73/32) & $+0.0001$ (58/25) \\
L14H26 & +0.40 & -- & $-0.0037$ (27/13) & $-0.0005$ (9/8)$^{\dagger}$ \\
L19H16 & -0.82 & $+0.0031$ (18/6)$^{\dagger}$ & $+0.0052$ (92/59) & $-0.0017$ (144/61) \\
L23H15 & -0.95 & -- & $-0.0118$ (30/26) & $-0.0009$ (32/25) \\
L26H23 & -0.37 & $+0.0000$ (1/1)$^{\dagger}$ & $+0.0092$ (18/16)$^{\dagger}$ & $+0.0033$ (39/22) \\
L13H10 & -0.68 & $+0.0152$ (2/2)$^{\dagger}$ & $-0.0157$ (8/5)$^{\dagger}$ & $+0.0028$ (31/19) \\
\bottomrule
\end{tabular}
\caption{Saved Stage-1 test-only RI events partitioned by the role of the \emph{current} token (first anchor): family-weighted raw target-minus-control gap, with events/families. $^{\dagger}$: below the 20-event/10-family support rule. Heads with no scored QK events in a partition are marked --. Source: \texttt{results/stage3\_inputs\_v1/current\_name\_diagnostics.json}.}
\label{tab:current}
\end{table}

For the positive movers with support in both partitions the name gap is clearly
negative when the current token is a control name and near zero otherwise: L18H19
$-0.023$ (111 events) vs.\ $+0.000$ (51); L21H18 $-0.017$ (194) vs.\ $+0.001$ (185);
L21H6 $-0.025$ (421) vs.\ $-0.004$ (116). For the anti-copying head L19H16 the control
partition is positive ($+0.005$, 92 events) and the neither partition slightly negative.
This is the pattern the copy-penalizing hypothesis predicts: a head whose OV circuit
raises its current token's logit scores badly on the name gap exactly at the events
where the current token is a control name. The partitions differ in position and
composition (control-token events are mostly self events on chain-middle names), so this
is an association within the saved events, not an isolated causal explanation of the
RI failure. Across the 33 heads with supported control-partition events the Spearman
correlation of the control-partition gap with $I_h$ is $-0.26$ ($p=0.15$); the sample is
too small for more than a direction. Many strong heads (L17H1, L27H6, L20H1, L30H18,
L15H25, L21H23, L17H17) have no or almost no scored QK events and are insufficient in
every partition.

\subsection{Attended-name mover classification}
\label{sec:attended}
\begin{figure}[ht]
\centering
\includegraphics[width=.75\textwidth]{fig9_attended_name.png}
\caption{Fraction of attended-name events (test-block rows of the core prompts whose
attention argmax lies inside a visible test name) in which the contextual output
$o^h_jW_U$ favours the attended name's first token over the other visible names,
against Scope-F importance. Dashed lines: the fixed 80\% thresholds.}
\label{fig:attended}
\end{figure}

Rule (fixed before results): a head is a \emph{name mover} if the contrast is positive
in $\ge80\%$ of its attended-name events and a \emph{negative name mover} if negative in
$\ge80\%$, with support $\ge20$ events in $\ge10$ families; below support the head is
``insufficient'', never ``neither''. Contrast = contextual logit of the attended name's
first token minus the mean over the other visible distinct names (both roles);
collisions of anchor token IDs are excluded. Sensitivity variants: last-token anchor,
same-role controls only, non-self events only, colon site only, base prompts only.
All 160 core prompts of all three replicas enter the classification (94{,}015
attended-name events over the 105 heads).

Results (Figure~\ref{fig:attended}): six name movers, L27H6 (96\% of 768 events, mean
contrast $+2.09$), L17H24 (94\% of 1{,}073), L26H31 (91\% of 621, mean $+2.09$),
L18H19 (84\%), L22H5 (82\%), L21H6 (82\%); one negative name mover, L20H1 (17\%
positive, mean $-0.10$). Six of the seven labels are consistent with the sign of
$I^F_h$; the seventh, L26H31, is a strong mover of the name it attends to with a
negligible individual causal effect ($I^F_h=+0.027$). Its attention rarely selects the
answer: 12.9\% of its 62 colon-site name events land on the answer (25\% of the 16
base-prompt colon events, nine families), and it puts 0.09 of its colon mass on the
query mother. Strong name movement with weak answer selection is a consistent
descriptive account of its small effect; it is not established as the explanation,
since redundancy with other movers and the specific intervention are alternatives.
L21H18 (79\%) and L19H22 (78\%) miss the threshold and are name movers on colon-site
events only; L18H18 (74\%) is not labelled. One support detail matters for the
strongest mover: L18H19's overall label rests on 646 events, but its colon-site subset
has only 17 events in nine families and is \emph{insufficient} under the frozen rule,
so the label describes its behaviour at earlier test positions, not specifically at
the causal answer position. Every label disappears at the last-token anchor: what these heads move is the
first token of the attended name, the token the model has to produce next.

The three pre-registered candidates give a mixed answer. L19H16 is negative in 77\% of
1{,}906 events (just under the rule; negative name mover on the colon site only) and
L26H23 in 79\% (negative name mover with same-role controls). L23H15, however, favours
the attended name in 64\% of 2{,}428 events (mean $+0.06$) while its causal effect is
$-0.95$: it is \emph{not} a negative name mover by its own output. Its attention is the
same as L27H6's (0.51 on the query mother, 100\% query-tracking), it is the strongest
induction/retrieval head of the inventory on the synthetic tasks (below), and its
direct vocabulary write at the colon is $+0.04$. The ``negative name mover''
description inherited from Stage 2 is therefore not supported for L23H15; how its
negative effect arises is open, and mediation by later components is the hypothesis to
test in Stage 4. L30H18 is different again: its contrast on the attended name is
positive (67\%, mean $+0.60$) --- it \emph{promotes} what it attends to --- and it
attends to the distractor. Promotion of the distractor name is a hypothesis consistent
with all three observations (attention, projection, causal sign); a targeted
intervention on its attention would be needed to establish it. The label ``negative
mover'' thus covers at least two different patterns; only L20H1 (and, marginally,
L19H16 and L26H23) suppresses the name it attends to.

L17H1 favours the attended name in 77\% of 3{,}469 events, but almost all are earlier-
site events; its colon-site support is insufficient, consistent with its zero effect
there.

\subsection{Matched contextual RI}
Replacing the raw current-token embedding by the actual normalized value input at the
same position, with event IDs, controls, anchors and family weighting held fixed
(70 heads with valid events; 27 inventory heads have no saved scored event and 8 have
too few matched events), does not repair the Stage-1 statistics. On the main
all-facts/all-test first-anchor name gap ($n=68$ supported heads) the Spearman
correlation with Scope-P importance moves from $-0.26$ (raw) to $-0.17$ (contextual);
with the RI-supported family/order subset of pairs, from $-0.28$ to $-0.19$. The
raw target statistic is uncorrelated ($+0.01$ raw, $+0.02$ contextual); the last-anchor
target statistic is negative under both ($-0.31$, $-0.23$). The embedding choice is
not what makes RI fail: RI scores the current position's own residual, whereas the
heads that matter act on what they attend to, or indirectly. This is a sensitivity
analysis on a selected inventory, not an all-head result.

\subsection{Attention--value factorial at the colon}
\begin{figure}[ht]
\centering
\includegraphics[width=\textwidth]{fig8_attention_value.png}
\caption{Ordinary Scope-F patch beside the two single replacements (corrupted attention
pattern with clean values; clean pattern with corrupted values) and their interaction,
178 pairs.}
\label{fig:av}
\end{figure}

For every head the value replacement alone carries the effect (L18H19: values $+3.16$,
pattern $+0.08$; L27H6: $+2.59$ vs.\ $+0.27$; L18H18 $+2.88$ vs.\ $+0.11$; L21H18
$+1.70$ vs.\ $+0.01$; L25H17 $+1.29$ vs.\ $-0.02$). The routing term is small and
non-zero only for the negative heads (L20H1 $-0.22$, L30H18 $-0.18$, L19H16 $-0.16$,
L17H3 $-0.10$, L23H15 $-0.09$) and for L27H6 ($+0.27$). The family-mean interaction
terms are also small ($-0.27$ for L18H18, $+0.21$ for L27H6, $-0.16$ for L18H19,
otherwise near zero), but these means cancel large family-level interactions of both
signs: the mean absolute interaction per family is 1.25 logits for L27H6 (maximum
4.58), 0.81 for L18H18 (maximum 3.27), 0.54 for L18H19 (maximum 3.25), 0.47 for L21H18
and 0.58 for L20H1, and the mean absolute routing term is 0.68 for L27H6 and 0.41 for
L18H18. Value replacement carries the larger effect on average for the leading
answer-position heads; a uniformly weak interaction between routing and values is not
shown, and separability should not be assumed in Stage 4. On average the picture
agrees with Section~\ref{sec:attention}: the corruption barely moves the family-mean
attention of these heads, so what changes is mostly the content at the attended
position and that content is what the reader carries to the colon. The factorial does
not separate ``name copying'' from ``reading a relational binding'': both predict value
dominance.

\subsection{Repeated-sequence and retrieval fingerprints}
\begin{figure}[ht]
\centering
\includegraphics[width=\textwidth]{fig7_synthetic.png}
\caption{Synthetic fingerprints for all 105 heads: attention to the copy source against
the head's output preference for the gold token, on the key--value retrieval task
(left) and the repeated random sequence (right).}
\label{fig:synth}
\end{figure}

The model solves both tasks (key--value retrieval 100\% of 32 trials; repeated
sequence 96.1\% of 992 probes), so both fingerprints are defined for every head.
Retrieval-type attention (argmax on the copy source) is shown by L19H16 (0.97 of
probes), L21H18 (0.88), L23H15 (0.84), L18H18 (0.56), L16H1 (0.47), L18H19 (0.19),
L22H5 (0.16), L26H23 (0.13); induction-type attention on the repeated sequence by
L23H15 (0.66), L21H6 (0.64), L16H1 (0.45), L21H18 (0.39), L19H16 (0.36), L22H5
(0.22), L26H23 (0.15). Random controls reach at most 0.44 and 0.17. Output sign on the
gold token separates the two kinds of source-attender: L21H6 ($+0.42$ on repeats),
L21H18 ($+0.15$/$+0.09$), L18H19, L22H5 and L16H1 promote the copied token; L19H16
($-0.24$/$-0.10$), L26H23 ($-0.19$/$-0.17$) and L23H15 ($-0.05$/$-0.02$) attend to
the source and \emph{do not} promote it. L25H17, L30H13 and L21H23 have positive gold
output with self-attention on both tasks, as in the kinship task. L27H6, L17H1, L18H18
(output), L20H1 and L30H18 show no fingerprint on either probe. Absence on these two
probes is not evidence of task specificity; other synthetic tasks might reveal a
transferable pattern. Transfer to a synthetic task supports a behavioural
capability and does not show that the same algorithm solves the kinship task.

\section{Head-profile table}
\label{sec:cards}
\begin{table}[p]
\centering\scriptsize
\setlength{\tabcolsep}{3.5pt}
\resizebox{\textwidth}{!}{%
\begin{tabular}{llrrlrrrrl}
\toprule
Head & Group & $I_h$ (P) & $I_h^{F}$ & Attention at colon & Ctx.\ out. & Copy & KV & Rep. & Attended-name label \\
 & & 178 pairs & 178 pairs & Q-mother/Q-child/D-mother/self & $\Delta$ mother & weights & src att. & src argmax & (fraction positive; events/families) \\
\midrule
L17H1 & S+ & $+5.70$ & $-0.06$ & 0.05/0.06/0.02/0.04 & $-0.00$ & $+0.01$ & 0.00 & 0.00 & neither (77\%; 3469/20) \\
L18H19 & S+RI & $+3.07$ & $+3.08$ & 0.13/0.05/0.03/0.00 & $+0.16$ & $+0.12$ & 0.16 & 0.02 & NM (84\%; 646/20) \\
L27H6 & S+ & $+3.07$ & $+3.06$ & 0.58/0.00/0.05/0.00 & $+2.99$ & $+0.04$ & 0.05 & 0.01 & NM (96\%; 768/20) \\
L18H18 & S+ & $+2.70$ & $+2.72$ & 0.14/0.03/0.03/0.00 & $+0.11$ & $-0.08$ & 0.27 & 0.06 & neither (73\%; 332/20) \\
L21H18 & S+ & $+1.70$ & $+1.70$ & 0.60/0.00/0.10/0.00 & $+0.25$ & $+0.11$ & 0.46 & 0.39 & neither (80\%; 1656/20) \\
L25H17 & S+ & $+1.27$ & $+1.27$ & 0.00/0.00/0.00/0.57 & $+0.15$ & $-0.01$ & 0.01 & 0.00 & neither (43\%; 625/20) \\
L22H5 & S+RI & $+0.98$ & $+0.98$ & 0.34/0.01/0.08/0.01 & $+0.41$ & $+0.04$ & 0.18 & 0.22 & NM (82\%; 730/20) \\
L16H21 & S+RI & $+0.92$ & $+0.89$ & 0.01/0.40/0.01/0.00 & $-0.01$ & $-0.02$ & 0.04 & 0.00 & neither (60\%; 551/20) \\
L16H1 & S+RI & $+0.87$ & $+0.86$ & 0.01/0.36/0.02/0.01 & $-0.00$ & $+0.01$ & 0.28 & 0.45 & neither (50\%; 669/20) \\
L21H6 & S+ & $+0.70$ & $+0.69$ & 0.29/0.01/0.04/0.00 & $+0.33$ & $+0.35$ & 0.12 & 0.64 & NM (81\%; 858/20) \\
L15H25 & S+ & $+0.68$ & $+0.17$ & 0.03/0.06/0.02/0.01 & $-0.00$ & $+0.04$ & 0.00 & 0.00 & neither (42\%; 1797/20) \\
L30H13 & S+ & $+0.63$ & $+0.63$ & 0.02/0.01/0.01/0.19 & $+0.76$ & $-0.03$ & 0.03 & 0.00 & neither (37\%; 917/20) \\
L21H23 & S+ & $+0.61$ & $+0.60$ & 0.05/0.01/0.02/0.20 & $+0.07$ & $-0.00$ & 0.04 & 0.00 & insuff. (25\%; 20/8) \\
L17H17 & S+ & $+0.53$ & $+0.51$ & 0.04/0.05/0.03/0.19 & $+0.00$ & $+0.01$ & 0.01 & 0.00 & neither (39\%; 103/19) \\
L17H24 & S+RI & $+0.43$ & $+0.43$ & 0.13/0.03/0.05/0.02 & $+0.03$ & $+0.18$ & 0.02 & 0.00 & NM (94\%; 1073/20) \\
L19H22 & S+ & $+0.43$ & $+0.42$ & 0.31/0.06/0.07/0.01 & $+0.07$ & $+0.17$ & 0.04 & 0.01 & neither (78\%; 1127/20) \\
L14H26 & S+ & $+0.40$ & $+0.01$ & 0.01/0.05/0.00/0.01 & $+0.01$ & $-0.00$ & 0.00 & 0.00 & neither (54\%; 959/20) \\
L26H23 & S$-$ & $-0.37$ & $-0.37$ & 0.21/0.02/0.03/0.01 & $-0.04$ & $-0.20$ & 0.16 & 0.15 & neither (21\%; 1670/20) \\
L18H24 & S$-$ & $-0.53$ & $-0.57$ & 0.06/0.06/0.04/0.02 & $+0.00$ & $-0.06$ & 0.02 & 0.00 & neither (34\%; 70/15) \\
L17H3 & S$-$ & $-0.55$ & $-0.40$ & 0.06/0.10/0.02/0.04 & $+0.00$ & $-0.03$ & 0.08 & 0.00 & neither (44\%; 319/20) \\
L13H10 & S$-$ & $-0.68$ & $-0.03$ & 0.00/0.01/0.00/0.42 & $+0.00$ & $-0.03$ & 0.00 & 0.00 & neither (43\%; 700/20) \\
L19H16 & S$-$ & $-0.82$ & $-0.83$ & 0.21/0.04/0.03/0.01 & $-0.04$ & $-0.09$ & 0.55 & 0.36 & neither (23\%; 1906/20) \\
L23H15 & S$-$ & $-0.95$ & $-0.95$ & 0.51/0.01/0.05/0.00 & $+0.04$ & $+0.09$ & 0.49 & 0.66 & neither (64\%; 2428/20) \\
L30H18 & S$-$ & $-1.06$ & $-1.06$ & 0.05/0.01/0.11/0.00 & $-0.88$ & $+0.60$ & 0.05 & 0.00 & neither (67\%; 2528/20) \\
L20H1 & S$-$ & $-1.56$ & $-1.53$ & 0.17/0.05/0.01/0.02 & $-0.09$ & $-0.17$ & 0.02 & 0.00 & neg.\ NM (17\%; 437/20) \\
\bottomrule
\end{tabular}%
}
\caption{Head-profile table for the 25 inventory heads with $|I_h|\ge 0.3$, ordered by Scope-P importance. Group codes: S+ strong positive outside RI; S+RI strong positive RI-selected; S$-$ strong negative. Attention: mean probability mass at the final colon on base prompts (20 common families) on the query-fact mother, the query-fact child, one distractor mother (mean of three) and the colon itself. Ctx.\ out.: contextual output projection $o^h_j W_U$ at the colon, first token of the query mother minus mean of the distractor mothers. Copy: weight copying score on name tokens (self minus mean of other names). KV / Rep.: synthetic retrieval and repeated-sequence fingerprints (attention to the copy source; argmax on the source). Attended-name label under the fixed 80\% rule (NM = name mover; Section~\ref{sec:attended}), with the fraction of events favouring the attended name and the event/family support.}
\label{tab:profiles}
\end{table}


Table~\ref{tab:profiles} summarizes the 25 strong heads; the per-head evidence cards
(\texttt{analysis/head\_cards.md}) and the full tables cover all 105. For every head
and analysis the status is one of completed, insufficient support, unavailable input,
or unresolved: contextual RI is insufficient for 35 heads (no or too few saved QK
events), the conditional synthetic scores are complete for all heads, the
attended-name label is insufficient for 19 heads (mostly heads whose attention never
lands on a name), and every other cell is completed.
No optional extension of Section~1.5 was run. The trigger for the controlled
representation readout is met for two sites --- the query-fact ``is''/period positions
of L17H1 and the child's last token for L15H25 are causally localized while the
vocabulary projection does not identify what they carry --- so a small readout on the 20
common families (intact vs.\ head-intervened residuals in the same readout prompt, with
no-injection, irrelevant-injection and entity-swap controls, probabilities recorded
alongside decoded text) is recommended; it needs a new small capture run, since the
saved exports do not contain the residual vectors. Broad spectral analysis and
function-vector mediation are not triggered.

\section{Evidence assessment of the Stage-2 hypotheses}
\label{sec:hyp}
\textbf{Copying.} Partly supported, with an important correction. Six heads
(L27H6, L17H24, L26H31, L18H19, L22H5, L21H6) promote the first token of the name they
attend to in $\ge80\%$ of events, and L27H6 writes the answer directly ($+2.99$
logits).
Weight copying scores have the right sign for these heads and for the anti-copiers
L20H1, L26H23, L19H16. The Stage-2 conjecture that the negative name-gap correlation
arises because ``the causal heads copy names'' is therefore consistent with the
positive movers. But copying is neither necessary nor sufficient: L18H18 (2.7 logits)
and L27H6 have weight scores near zero or negative, the strongest weight copier L30H18
hurts the answer by copying the distractor, L26H31 moves names as strongly as L27H6 but
attends to the wrong ones and has no effect, and the largest head, L17H1, does not act
at the colon at all.

\textbf{Contextual information.} Supported for the location, not yet for the content.
The heads' effects come mainly through their values (factorial, on average), their
attention tracks the queried fact under query change (and on average under relocation),
and their raw-embedding projections
are uninformative while the contextual projections are partly informative. Contextual
RI at the current position does not help, because the information the readers use is
at the attended position, written there by earlier layers (Section~\ref{sec:position}).
What that information encodes (name identity, a relation binding, or both) is a
Stage-4 question.

\textbf{Early vs.\ late positions.} Supported and sharpened. The effect of L17H1,
L15H25, L14H26 and L13H10 is written inside the query-fact sentence (L17H1 at the
mother tokens, ``is'' and the period; L15H25 at the child's last token) and read
later; all other strong heads act only at the colon. The earlier ``layer-6--11
writers'' of Stage 2 (L6H24, L7H1, L8H15, L9H16, L12H17) are confirmed to have zero
colon effect, but their all-position effects are small ($\le0.28$); the important
writer is L17H1, in the same layer band as the readers.

\textbf{What Stage 3 does not establish.} No edge between any writer and any reader has
been tested; single-position additivity within one head says nothing about additivity
across heads; the labels ``name mover'' and ``negative name mover'' describe the sign of
one head's output on its attended token and are not mechanism claims; the changed-query
control is a discovery measurement on 20 families; and all thresholds are descriptive.

\section{Hand-off to Stage 4}
\label{sec:stage4}
Stage 4 receives the following components and specific communication hypotheses.

\begin{enumerate}
\item \textbf{Writer L17H1 $\rightarrow$ mother-attending readers.} L17H1 writes at
the query mother's tokens, ``is'' and the period; L27H6, L21H18, L22H5, L18H19,
L18H18 attend to the mother's tokens at the colon. Path patching L17H1's output at
those positions into the readers' keys vs.\ values will test whether L17H1 supplies the
content the readers carry (values) or the address they select (keys). The value
dominance of the factorial predicts values.
\item \textbf{L15H25 $\rightarrow$ child-attenders.} L15H25 writes at the query child's
last token; L16H21 and L16H1 attend to the child at the colon and have no direct
vocabulary output. Test whether they route through L15H25's write and where their
$0.9$-logit effects surface (which later heads read them).
\item \textbf{Heads whose projection does not explain their effect.} L18H19, L18H18,
L20H1, L23H15, L19H16 have colon effects far larger than their vocabulary projections.
Test the mediation hypothesis: path-patch their output into later heads and MLPs at the
colon, and compute a matched clean-minus-donor projection under an explicit
normalization reference; L23H15's and L30H18's negative effects in particular need this
before any type label. L17H3 ($F/P=0.74$) is a mixed-position candidate whose earlier
sites should be located first.
\item \textbf{Positive vs.\ negative movers.} Joint interventions on \{L27H6, L21H18,
L22H5, L21H6\} against \{L20H1, L19H16, L26H23\} to test redundancy and whether the
negative heads calibrate the positive ones; group ablation before any faithfulness
claim.
\item \textbf{Validation.} Only after edges are established: minimal-circuit
faithfulness on the 87 held-out families and the reorder/changed-query conditions.
\end{enumerate}

\section{Conclusion}
Stage 3 replaces the Stage-2 sketch with a measured profile of every inventory head.
The largest causal head, L17H1, is not a reader at the answer position: it acts inside
the query-fact sentence, at and after the mother's name, and its effect is recovered
symmetrically by reverse patching. Twenty strong heads deliver their whole effect at
the colon, mainly through their values, with attention that follows the queried fact
under a changed query and, on average, under relocation and name swapping. Among them,
six promote the first token of the name they attend to (L27H6 directly writing the
answer at $+3$ logits; L26H31 as strongly, but at the wrong names and without effect),
one suppresses it (L20H1), several have effects that their own vocabulary projection
does not explain (L18H19, L18H18, L23H15, L19H16), and one attends to the distractor
and promotes it, consistent with its negative effect (L30H18). Neither the raw-embedding RI projection nor its
contextual variant tracks these effects, because the relevant information sits at the
attended positions and, for the indirect heads, downstream. Stage 3 therefore ends
with explicit components and edges to test, not with a circuit.

\end{document}
